\chapter{Problem Statement}


\begin{itemize}
    
\end{itemize}

% ----------------------------------------------------------------
\section{Background and Motivation}
% ----------------------------------------------------------------

Traffic congestion in major urban areas has become one of the most pressing
challenges in modern urban infrastructure management. In Vietnam, Ho Chi Minh City
suffers estimated losses of approximately \$6 billion USD annually~\cite{hcmc_congestion_2022},
while Hanoi loses around \$1.2 billion USD per year due to traffic
jams---equivalent to nearly 17\% of the city's logistics GDP~\cite{hanoi_congestion_worldbank}.
These losses reflect an accumulation of indirect costs: wasted labor productivity,
excess fuel consumption by idling vehicles, and healthcare costs associated with
vehicle emissions.

A root cause of this problem lies in the inefficiency of traditional traffic signal
control systems. Fixed-time signal control operates on pre-programmed cycles with
no ability to adapt to real-time fluctuations in traffic
demand~\cite{adaptive_signal_survey_2024}. As a result, during peak hours, vehicle
queues at intersections routinely exceed the capacity of the current green phase,
while other phases waste green time on near-empty lanes---squandering the
throughput capacity of the entire network.

To address this limitation, machine learning approaches---particularly
\textit{Reinforcement Learning}---have emerged as a promising solution for
adaptive traffic signal control~\cite{rl_tsc_survey_springer2026,
marl_tsc_sustainability2023}. Unlike classical optimization methods, RL does not
require an explicit mathematical model of the environment: an agent learns improved
decision-making through repeated interaction with a simulated environment, receiving
rewards based on achieved traffic performance. When the problem scales to
multi-intersection networks, \textit{Multi-Agent Reinforcement Learning}
becomes a natural approach, where each intersection is controlled by an independent
agent while sharing knowledge through a parameter-sharing
mechanism~\cite{drl_tsc_survey_sciencedirect2024}.

The motivation for this research stems from the practical urgency of congestion in
Vietnamese cities: intelligent traffic signal control offers a deployable solution
that does not require large-scale physical infrastructure investment. The use of
real traffic demand data from Cologne, Germany, via the TAPAS Cologne
dataset~\cite{tapas_cologne_sumo}---recognized as one of the largest freely
available traffic simulation datasets---enables evaluation of RL algorithms under
complex, realistic urban traffic conditions that go well beyond synthetic benchmarks.

% ----------------------------------------------------------------
\section{Problem Description}
% ----------------------------------------------------------------

The problem addressed in this thesis is \textbf{optimizing traffic signal control
over a real-world urban road network using multi-agent reinforcement learning}.
Specifically:

\begin{itemize}
    \item \textbf{Input:} Local traffic state at each intersection at every time
    step, including vehicle counts, queue lengths, and current signal phase,
    collected via the TraCI interface of the SUMO traffic simulator.

    \item \textbf{Output:} Signal phase control actions (selection of the next
    phase) produced by each agent at its assigned intersection, updated
    continuously in the simulation environment.

    \item \textbf{Scope:} Two real-world traffic scenarios from Cologne,
    Germany---Cologne-3 and Cologne-8 ---built on the TAPAS Cologne dataset. Three RL
    algorithms are trained and evaluated: PPO, DQN, and DDQN, all under an
    Independent architecture with parameter sharing (IPPO, IDQN, IDDQN). Three
    reward functions---\texttt{queue}, \texttt{pressure}, and
    \texttt{wait-clip}---are investigated to analyze the effect of reward design
    on learned control policies.

    \item \textbf{Assumptions:} Each agent observes only local information at its
    own intersection; no explicit communication occurs between agents during
    decision-making. The SUMO simulation environment is assumed to be a
    sufficiently representative proxy for real-world urban traffic conditions
    within the scope of this study.
\end{itemize}

% ----------------------------------------------------------------
\section{Research Objectives and Scope}
% ----------------------------------------------------------------

The overall objective of this thesis is to \textit{implement, train, and
comprehensively evaluate three MARL algorithms (PPO, DQN, DDQN) for traffic signal
control on real-world urban road networks, benchmarked against conventional
fixed-time control}. To achieve this, the following specific tasks are addressed:

\begin{enumerate}
    \item Construct a traffic simulation environment in SUMO and formulate the
    problem as a Markov Decision Process, defining the state space, action
    space, and reward function.

    \item Implement three RL algorithms (IPPO, IDQN, IDDQN) under an Independent
    architecture with parameter sharing across the Cologne-3 and Cologne-8
    scenarios.

    \item Investigate the effect of three reward functions (\texttt{queue},
    \texttt{pressure}, \texttt{wait-clip}) on learned control policies and
    multi-agent coordination behavior.

    \item Conduct a comprehensive performance comparison among the MARL algorithms
    and against fixed-time control, using standard traffic metrics ( queue length, and throughput).
\end{enumerate}

The main contributions of this thesis are: (i) successful training and deployment
of three RL algorithms under the Independent architecture; (ii) quantitative
analysis of the effect of reward function design on multi-agent coordination;
and (iii) identification of the most
effective control configuration for the multi-intersection environment.

% ----------------------------------------------------------------
\section{Thesis Structure}
% ----------------------------------------------------------------

The remainder of this thesis is organized as follows:

\begin{itemize}
    \item \textbf{Chapter 2} presents the theoretical background on reinforcement
    learning, including MDP concepts, Q-learning, and detailed descriptions of the
    PPO, DQN, and DDQN algorithms.

    \item \textbf{Chapter 3} describes the system design, covering the SUMO
    simulation environment, state and action representations, reward functions and
    model configurations

    \item \textbf{Chapter 4} presents experimental results, including per-scenario
    specialist performance and reward function analysis.

    \item \textbf{Chapter 5} provides concluding remarks and directions for future
    research.
\end{itemize}